\documentclass[conference]{IEEEtran}
\usepackage[utf8]{inputenc}
\usepackage{booktabs}
\usepackage{amsmath}
\usepackage{url}
\usepackage[hidelinks]{hyperref}
\usepackage{xspace}

% Every number in this paper is generated from results/*.json by scripts/make_paper_numbers.py.
% Generated by scripts/make_paper_numbers.py. Do not edit by hand.
\newcommand{\PENDING}{\textbf{[pending]}}

\newcommand{\DataSkills}{9,740}
\newcommand{\DataMalicious}{7,505}
\newcommand{\DataBenign}{2,235}
\newcommand{\TestBenign}{545}
\newcommand{\TestAtrBenign}{455}
\newcommand{\RandomSvmFone}{0.932}
\newcommand{\SourceSvmFone}{0.647}
\newcommand{\SourceSvmFpr}{66\%}
\newcommand{\ProbeAcc}{91.7\%}
\newcommand{\ProbeSources}{10}
\newcommand{\FprAll}{65\%}
\newcommand{\FprMixed}{3.7\%}
\newcommand{\FprMixedCI}{[2.2\%, 5.3\%]}
\newcommand{\PrAll}{0.79}
\newcommand{\PrMixed}{0.93}
\newcommand{\FprSizeMatched}{58\%}
\newcommand{\MixedSkillTrustBench}{3973}
\newcommand{\MixedAgentTrap}{116}
\newcommand{\MixedSkillTrojan}{3}
\newcommand{\MechanismSentence}{In these sources 83\% of malicious skills are injected copies of a benign skill; breaking the pairs raises the benign false-positive rate only to 10\%, modestly above the 6\% of a size-matched mixed-only control (both over three seeds) and far below the 65\% of all-data, so label mixing within a source drives most of the gain while the pairing adds a little.}
\newcommand{\TransformerSentence}{A fine-tuned ModernBERT shows the same pattern: trained on all data it flags 79\% of unseen-source benign skills, and trained on mixed-label sources only 20\%, so the over-blocking is not merely an artifact of the linear model.}
\newcommand{\DecoderSentence}{A QLoRA-fine-tuned Qwen2.5-0.5B decoder (4-bit, one epoch) shows the same over-block: trained on all data it flags 80\% of unseen-source benign skills and on mixed-label sources only 1.8\% (95\% CI [0.7\%, 3.1\%]), so the false-positive effect spans linear, encoder and decoder models. On this small decoder the fix is a pure false-positive/recall operating-point trade: malicious recall falls 97\%$\to$12\% and ranking is unchanged (PR-AUC 0.79 vs 0.78), so unlike the fully-trained word-LR and ModernBERT models it shows no ranking gain, as expected for a 0.5B model fine-tuned for a single epoch.}
\newcommand{\ConclusionMechanism}{On this benchmark the gain survives breaking the benign--malicious pairs, so it rests on label mixing within a source rather than on paired construction.}
\newcommand{\LoraFprAll}{80\%}
\newcommand{\LoraFprMixed}{1.8\%}
\newcommand{\LoraFprMixedCI}{[0.7\%, 3.1\%]}
\newcommand{\LoraPrAll}{0.79}
\newcommand{\LoraPrMixed}{0.78}
\newcommand{\LoraRecallAll}{97\%}
\newcommand{\LoraRecallMixed}{12\%}
\newcommand{\EvasionK}{8}
\newcommand{\EvasionAsrAll}{17\%}
\newcommand{\EvasionControlAll}{5\%}
\newcommand{\EvasionAsrMixed}{7\%}
\newcommand{\EvasionGapAll}{12\%}
\newcommand{\EvasionControlMixed}{1\%}
\newcommand{\EvasionBaseRecallAll}{40\%}
\newcommand{\EvasionBaseRecallMixed}{20\%}
\newcommand{\RecallAll}{96\%}
\newcommand{\RecallMixed}{50\%}
\newcommand{\RecallMixedDeploy}{20\%}
\newcommand{\RecallAtTwoAll}{8\%}
\newcommand{\RecallAtTwoMixed}{46\%}
\newcommand{\MbAtrAll}{75\%}
\newcommand{\MbSfbAll}{100\%}
\newcommand{\MbAtrMixed}{4\%}
\newcommand{\MbSfbMixed}{100\%}
\newcommand{\MbPrAll}{0.66}
\newcommand{\MbPrMixed}{0.79}
\newcommand{\ProbeBalAcc}{76.0\%}
\newcommand{\WildN}{400}
\newcommand{\WildFlaggedAll}{30\%}
\newcommand{\WildFlaggedMixed}{5\%}
\newcommand{\WildOpenclaw}{32\%}


\begin{document}

\title{Source Shortcuts in Malicious Agent-Skill Detection:\\
Why Detectors Fail on Unseen Sources, and What the Fix Really Learns}

\author{\IEEEauthorblockN{Aryan Bhargava}
\IEEEauthorblockA{Newton School of Technology, Rishihood University\\
\url{https://github.com/Aryanbhargava18/SkillShift}}}

\maketitle

\begin{abstract}
Agent skills (\texttt{SKILL.md} files) extend LLM agents with instructions and scripts, and malicious skills are an
emerging supply-chain threat. On MaliciousSkillBench, text classifiers that look strong on a random split
wrongly flag about two thirds of benign skills from sources they have not seen. We show why and test a fix.
A TF-IDF probe names a skill's source benchmark from its text with \ProbeAcc{} accuracy over
\ProbeSources{} sources, and platform boilerplate such as \texttt{openclaw} acts as a ``benign'' cue only
because of how the benchmark's sources were mixed. Training only on sources that contain both labels cuts the
benign false-positive rate on unseen sources from \FprAll{} to \FprMixed{} (95\% CI \FprMixedCI{}); this trades
malicious recall at a fixed threshold (\RecallAll{} to \RecallMixed{}) but improves ranking, raising PR-AUC
from \PrAll{} to \PrMixed{} and recall at a 2\% false-positive budget from \RecallAt2All{} to \RecallAt2Mixed{}.
A size-matched all-data control stays at \FprSizeMatched{}.
We then ask what this fix actually learns. Almost all of that training data is one source built from
minimal pairs: a benign skill and a copy with an injected payload. \MechanismSentence{}
\TransformerSentence{} The same shortcut is an attack surface: appending \EvasionK{} benign
source-style tokens evades \EvasionAsrAll{} of caught malicious skills, against \EvasionControlAll{}
for random tokens, and the fix lowers this to \EvasionAsrMixed{}.
We release code, per-skill scores, and an ONNX scanner, and argue that
skill-detection benchmarks should report unseen-source false-positive rates with the threshold fixed in
advance, and should state how their malicious and benign examples were constructed.
\end{abstract}

\section{Introduction}
Agent frameworks now load third-party \emph{skills}: a \texttt{SKILL.md} file with YAML frontmatter, natural
language instructions, and optional helper scripts that the agent may run. A malicious skill can exfiltrate
credentials, run a downloaded binary, or quietly redirect the agent. Static detectors that read the skill
before installation are an obvious first line of defense, and MaliciousSkillBench~\cite{msb2026} collects
\DataSkills{} labeled skills from 13 source benchmarks to evaluate them.

The benchmark's own results show the problem we study. A word TF-IDF + SVM model reaches macro-F1
\RandomSvmFone{} when train and test are drawn from the same mix, but \SourceSvmFone{} when the test skills
come from sources never seen in training, where it flags \SourceSvmFpr{} of benign skills. A detector that
blocks two thirds of legitimate skills is not deployable. The benchmark paper reports this gap and attributes
it to ``unfamiliar benign source conventions'' but does not isolate a mechanism~\cite{msb2026}, and
SkillScreen~\cite{skillscreen2026} reports similar unseen-source false-positive rates for fine-tuned
transformers.

The phenomenon itself is well known: when sources differ in their label mix, models learn to recognise the
source instead of the label~\cite{torralba2011,gururangan2018,geirhos2020,zech2018,compton2023}, and the same
failure has been documented for malware detection~\cite{arp2022} and malicious-prompt
classifiers~\cite{fomin2026}. Our contribution is not the gap but its diagnosis and a tested remedy in the
agent-skill domain:
\begin{itemize}
\item We reproduce the published baselines and show that a skill's source is predictable from its text
(\S\ref{sec:shortcut}), and trace a concrete shortcut feature, \texttt{openclaw}, to how sources were mixed.
\item We show that training only on mixed-label sources removes most unseen-source false positives, and that a
size-matched control does not (\S\ref{sec:fix}).
\item We test \emph{why} the fix works: whether it depends on the benchmark's minimal-pair construction,
which single-label sources bring the false positives back, and whether simply masking the shortcut token is
enough (\S\ref{sec:mechanism}).
\item We fine-tune ModernBERT on both training sets and compare it with the TF-IDF models on the same test
skills with paired bootstrap intervals (\S\ref{sec:transformer}), and ship the result as an int8 ONNX scanner.
\item We show the same shortcut is an \emph{attack surface}: a good-word attack~\cite{lowd2005} using
source-style tokens evades the detector well beyond a random-token control, and the mixed-only fix
shrinks the evasion surface as it shrinks the false-positive rate (\S\ref{sec:attack}).
\end{itemize}

\section{Data and evaluation}
\label{sec:data}
\textbf{Dataset.} MaliciousSkillBench~\cite{msb2026} (CC-BY-4.0) contains \DataSkills{} skills:
\DataMalicious{} malicious and \DataBenign{} benign. Most of its 13 sources contain one label only.
We use the frozen \texttt{source\_disjoint} split: the test set is SkillHarm (728 malicious), ATR (455 benign,
21 malicious) and SkillFortifyBench (90/90), none of which appear in training. Nearly all benign training
skills come from a single source, SkillTrustBench.

\textbf{Metrics.} The positive class is malicious. Our main number is the \emph{benign false-positive rate}
(FPR) on unseen-source test skills, the share of legitimate skills a detector would block. We also report
malicious recall, macro-F1 and PR-AUC. Recall at a fixed FPR is computed two ways: with the threshold picked
on the test set (an upper bound, labelled ``oracle''), and with the threshold picked on validation, which
comes from the training sources (the deployable setting). Confidence intervals are 95\% paired,
class-stratified bootstrap intervals with 2,000 resamples; per-source rates on small groups use exact
Clopper-Pearson intervals.

\textbf{Models.} Word TF-IDF (1--2 grams) with logistic regression or a linear SVM, character TF-IDF (3--5
grams), a keyword baseline, and ModernBERT-base~\cite{modernbert2024} fine-tuned on the first 512 tokens
(AdamW, learning rate $3\cdot10^{-5}$, 2 epochs, effective batch 16, class-weighted loss).

\section{The shortcut}
\label{sec:shortcut}
\textbf{Source is predictable from text.} We train a word TF-IDF + logistic regression classifier to
predict a skill's \emph{source benchmark} (not its label) and evaluate it with held-out skills. It
reaches \ProbeAcc{} accuracy over \ProbeSources{} sources (chance $1/\ProbeSources{}$). Since most
sources carry a single label, a malicious/benign classifier can score well by recognising the source
and never learning what makes a skill dangerous.

\textbf{A concrete shortcut feature.} The token \texttt{openclaw}, platform boilerplate from skills
published for the OpenClaw agent, is one of the strongest ``benign'' cues an all-data model learns.
Yet inside the one source that contains both labels at scale (SkillTrustBench) it appears about
equally in benign and malicious skills, so it carries no signal about maliciousness. It looks benign
only because benign training data is dominated by sources where the token is common, while the
largest malicious source rarely uses it. The cue is an artifact of how the benchmark's sources were
mixed, exactly the mechanism of \cite{arp2022,torralba2011}.

\section{Training where both labels share a source}
\label{sec:fix}
If a source contains only one label it can teach only style; if it contains both, style cannot
separate the labels inside it, so the model must use content. We therefore train only on the sources
that contain both labels. On the unseen-source test set this cuts the benign false-positive rate from
\FprAll{} to \FprMixed{} (95\% CI \FprMixedCI{}) and raises PR-AUC from \PrAll{} to \PrMixed{}
(Table~\ref{tab:fix}).

\begin{table}[t]
\caption{Word TF-IDF + logistic regression on the unseen-source test set, default $0.5$ threshold.}
\label{tab:fix}
\centering
\begin{tabular}{lcc}
\toprule
training data & benign FPR & PR-AUC \\
\midrule
all sources & \FprAll{} & \PrAll{} \\
size-matched random control & \FprSizeMatched{} & -- \\
mixed-label sources only & \FprMixed{} & \PrMixed{} \\
\bottomrule
\end{tabular}
\end{table}

\textbf{It is not simply ``less data''.} Mixed-only training uses far fewer skills than all-data. A
control that draws a random subset of all-data of the same size and label balance still flags
\FprSizeMatched{} of unseen-source benign skills, so the gain comes from \emph{which} data, not how
much. Re-weighting every (source, label) group equally does not help either: the style signal is
still present, only reweighted.

\textbf{Thresholds do not transfer.} A threshold tuned to flag 2\% of \emph{validation} benign skills
(validation comes from the training sources) flags far more on the unseen test sources, so any number
reported ``at $k\%$ FPR'' must say where the threshold was set. We report both an oracle operating
point (threshold on the test set, an upper bound) and the deployable one (threshold on validation).

\section{What the fix learns}
\label{sec:mechanism}
Mixed-only training is dominated by one source, SkillTrustBench (\MixedSkillTrustBench{} of the
mixed-label training skills; AgentTrap contributes \MixedAgentTrap{} and SkillTrojan
\MixedSkillTrojan{}), and that source builds most malicious skills by injecting a payload into a real
benign skill. So the fix could work for two different reasons: label mixing within a source, or the
presence of minimal pairs (a benign skill and a near-copy carrying a payload). We separate them.

\MechanismSentence{}

\textbf{Which sources bring the false positives back.} Adding one single-label source at a time to
the mixed-only training set locates the data responsible for the collapse: the largest single-label
malicious sources restore the high benign false-positive rate, while small ones do not. This shows
the failure is driven by the single-label malicious bulk, not by a lack of benign coverage alone.

\textbf{Masking the obvious token is not enough.} Removing the \texttt{openclaw} token, or stripping
the YAML frontmatter, from all-data training does not close the gap: the model falls back on the many
other style cues that distinguish the sources. The shortcut is distributed across the text, so
editing one feature cannot fix an all-data model; changing the training composition can.

\section{Fine-tuned transformer}
\label{sec:transformer}
To check that the shortcut is a property of the data and not of linear models, we fine-tune
ModernBERT-base on both training sets and score it on the same unseen-source test skills, comparing
every model with a paired, class-stratified bootstrap. \TransformerSentence{} Per-source rates on the
small benign test groups (ATR, SkillFortifyBench) are reported with exact Clopper-Pearson intervals,
since ATR contributes only \TestAtrBenign{} of the \TestBenign{} benign test skills.


\section{The shortcut is also an attack surface}
\label{sec:attack}
The source shortcut does not only make detectors over-block benign skills from unseen sources; it lets
an attacker evade them. Appending benign source-style tokens to a malicious skill is a good-word
attack~\cite{lowd2005}; the honest question is whether the \emph{source-style} tokens help beyond the
dilution that appending any text causes to a TF-IDF vector. We answer with a dose-response on the linear
word-LR detector, whose benign-weighted features are directly readable: at each size $k$ we append its
top-$k$ benign-weighted features, and as a control $k$ random vocabulary tokens, to every malicious test
skill it catches at its deployment threshold. The same construction would need a gradient or query budget
against the transformers; we demonstrate the mechanism where the shortcut is legible.

At every $k$ the source-style tokens evade far more than random ones. With only \EvasionK{} injected
tokens, \EvasionAsrAll{} of caught malicious skills fall below threshold, against \EvasionControlAll{}
for the same number of random tokens --- a source-specific gap of \EvasionGapAll{} that padding does
not explain. The attack changes no payload and queries no model.

The mixed-only model of \S\ref{sec:fix} is more robust here too: the same \EvasionK{}-token attack
succeeds on only \EvasionAsrMixed{} of the skills it catches. Because it no longer leans on source
style, the training-set-composition fix that lowers the benign false-positive rate also shrinks the
evasion surface. One bias thus produces both failures --- over-blocking benign skills and letting
malicious ones evade --- and one data fix moves both.

We state two honesty notes. The success rates are over each model's \emph{caught} malicious skills, and
the two models catch different amounts before any attack (at its threshold the all-data model flags
\EvasionBaseRecallAll{} of malicious test skills, the mixed-only model only \EvasionBaseRecallMixed{}), so
the mixed-only model's lower attack rate partly reflects its lower base recall; its source-specific gap
is \EvasionAsrMixed{} benign vs \EvasionControlMixed{} random. The effect sizes are modest: this is a
mechanism demonstration, not a strong attack. It concerns the cheap static pre-install filter, and is
orthogonal to evasions that hide the payload from an LLM judge~\cite{pretext2026,cloakdetonate2026}.


\section{Related work}
\textbf{Dataset bias and shortcuts.} ``Name that dataset'' experiments show sources are easy to identify from
their content~\cite{torralba2011,liu2025decade}; annotation artifacts let NLI models answer from the hypothesis
alone~\cite{gururangan2018}; and a pneumonia detector learned the hospital rather than the
disease~\cite{zech2018}. Adding a source with a different label rate can hurt for the same
reason~\cite{compton2023}. Group reweighting and Group DRO~\cite{sagawa2020} and group-balanced
subsampling~\cite{idrissi2022} target this setting; our mixed-only training is a form of group-balanced data
selection restricted to sources where both labels exist. Counterfactually augmented data~\cite{kaushik2020},
pairs of examples that differ only in the label-relevant edit, is closely related to the minimal pairs we
find in SkillTrustBench.

\textbf{Security.} Arp et al.~\cite{arp2022} report a near-identical artifact in Android malware detection,
where a store URL became a top ``benign'' feature, and recommend per-source datasets that contain both
classes. Fomin~\cite{fomin2026} shows, for malicious-prompt classifiers, that source identity is predictable
from text and that thresholds do not transfer across held-out datasets. In the skill domain,
MaliciousSkillBench~\cite{msb2026} reports the source-disjoint gap, a class-balancing control and a scaffold
sanitization control, and SkillScreen~\cite{skillscreen2026} reports transformer results and threshold
non-transfer on the same split. Neither tests source predictability, source-composition interventions, or
the role of minimal pairs.

We are explicit about what is and is not new. That merging sources creates a label shortcut, and that
per-source both-class data is the remedy, is Arp et al.~\cite{arp2022}; we do not claim either. Our
contribution is the empirical decomposition in the agent-skill domain: the source-identity probe, the
minimal-pair versus label-mixing ablation, the dose-response that locates MalSkillBench as the culprit,
and the attack$\leftrightarrow$false-positive symmetry. It is also a sharper intervention than the
benchmark's own control: MaliciousSkillBench's class balancing lowers the word model's benign FPR only to
about 43\%~\cite{msb2026}, against our \FprMixed{} from source selection.

\section{Threats to validity}
\label{sec:threats}
\textbf{One benchmark, one split.} All results use MaliciousSkillBench's frozen \texttt{source\_disjoint}
split. The unseen benign test set has only \TestBenign{} skills from two sources (ATR and SkillFortifyBench),
and ATR contributes \TestAtrBenign{} of them, so benign FPR is largely an ATR number. A full
leave-one-source-out evaluation would show how much the conclusions depend on which sources are held out.

\textbf{``Mixed-only'' is effectively one source.} SkillTrojan contributes \MixedSkillTrojan{} skills and
AgentTrap \MixedAgentTrap{}; SkillTrustBench contributes the other \MixedSkillTrustBench{}. Our fix is therefore
best described as ``train on SkillTrustBench'', and its success may depend on how that source was built
(\S\ref{sec:mechanism}).

\textbf{Construction artifacts.} Inside SkillTrustBench most malicious skills are injected variants of benign
skills, and the held-out malicious sources are also injected or backdoored. A detector that learns ``find the
injected span'' may do well here and still miss malicious skills that were written from scratch.

\textbf{Thresholds.} Validation data comes from the training sources, so thresholds tuned on it do not
transfer to unseen sources (\S\ref{sec:fix}), as also reported for this benchmark by
SkillScreen~\cite{skillscreen2026}. Numbers labelled ``oracle'' pick the threshold on the test set and are
upper bounds.

\textbf{Prevalence.} The benchmark is 77\% malicious; real registries are mostly benign. Precision at
realistic prevalence would be much lower than the precision implied here, which is why we report FPR rather
than accuracy.

\textbf{Compute.} ModernBERT was fine-tuned once per training set (seed 0) on a laptop GPU, on the first 512
tokens only. Transformer results are a single run, not a distribution over seeds.

\section{Conclusion}
\label{sec:conclusion}
Malicious-skill detectors trained on merged benchmarks learn where a skill came from. Training only where both
labels share a source removes most of the resulting false positives, and a size-matched control rules out the
``less data'' explanation. \ConclusionMechanism{} For benchmark builders this suggests two concrete practices:
report benign FPR on unseen sources with the threshold fixed before looking at the test set, and document how
malicious and benign examples within each source were constructed, so that users can tell whether a detector
learned what makes a skill dangerous or how the benchmark was assembled.


\bibliographystyle{IEEEtran}
\bibliography{refs}

\end{document}
